% 5.1 owner: Eileen (protocol). 5.2-5.4 owner: Prisca.
\section{Method}
\label{sec:method}

Every experiment follows the same pipeline (Figure~\ref{fig:pipeline}). The model first judges the original
program, then judges a changed version in which the verifier confirms the vulnerability is still present, so a
\safe{} verdict on it is a genuine miss.

\begin{figure}[t]
\centering
\begin{tikzpicture}[font=\scriptsize, node distance=0.35cm,
  box/.style={draw, rounded corners=2pt, align=center, minimum height=0.9cm, text width=1.45cm},
  arr/.style={-{Stealth[length=1.6mm]}}]
\node[box] (b) {\textbf{1. Baseline}\\ model judges the original};
\node[box, right=of b] (c) {\textbf{2. Change}\\ QA edit, tarpit, app, or attack};
\node[box, right=of c] (r) {\textbf{3. Re-judge}\\ same question, changed code};
\node[box, right=of r, draw=red!70!black, thick] (v) {\textbf{4. Verify}\\ run it with attack inputs};
\draw[arr] (b) -- (c); \draw[arr] (c) -- (r); \draw[arr] (r) -- (v);
\end{tikzpicture}
\caption{Evaluation pipeline. Step 4 executes every changed program, so we never assume the bug survived the change.}
\label{fig:pipeline}
\end{figure}

% Owner: Eileen. Replace with your text (keep the \subsection and the label).
% Outline (section 5): model comparison (only gpt-oss:20b recognizes safe code: 32/36 vs <=14/36 on the 24-stub
% sample; Qwen2.5-Coder 7B specificity 6%); frozen prompt (tag prompts-v1); sampling per publisher (gpt-oss:
% temperature 1.0, top_p 1.0, top_k off, repeat penalty off), 3 repeats, seed 17739 + repeat (Part 1: Qwen's
% settings, 1 repeat); Part 1 pipeline (baseline -> QA edit -> re-evaluation in the same and a fresh conversation ->
% verifier; outcomes held / genuine miss / incidental fix / fixed-but-flagged; misleading edits as an upper-bound
% probe); Part 4 attacks (fake end of the code block, unseen helper, instruction comment, instruction in a string,
% fake prior verdict, fake tests, pushback with pressure / authority, comment moved into app.js); metrics (miss rate,
% paired misses, false-alarm rate, 95% Wilson intervals, exact McNemar paired by program and repeat).
\subsection{Models, Prompts and Evaluation Protocol}
\label{sec:protocol}
\todo{Eileen: models, prompts and sampling, Part 1 pipeline, Part 4 attacks, metrics.}
   % Eileen: models, prompts, sampling, Part 1, Part 4 attacks, metrics

\subsection{Part 2: Testability tarpit transformations}
\label{sec:tarpits}
Al Kassar et al.\ identified code patterns, ``testability tarpits'', that make static and dynamic security testing
tools miss vulnerabilities in web applications~\cite{alkassar2022tarpits}. Rui Wang (CMU) shared a list of 80
JavaScript transformations that introduce such patterns. We re-implemented the nine he ranked most natural, by
how common they are in real code and how small a change they make: \emph{arrow-function},
\emph{computed-properties}, \emph{nested-function}, \emph{nullish-coalescing}, \emph{object-create},
\emph{simple-array}, \emph{ternary}, \emph{try-catch} and \emph{weakref}.

Our generator (\texttt{transforms/tarpit.js}) parses each stub with \texttt{acorn} and rewrites only the expression
that reaches the sink. For example, the ternary transformation turns \texttt{exec(cmd)} into
\texttt{var \_x\_0 = 5; exec(\_x\_0 > 9 ? 'safe' : cmd)}, which always runs the original command. Edits are
spliced into the original text, so comments and formatting are kept. Helper statements are inserted just before
the statement that holds the sink, and expression bodies are first wrapped in braces, so nothing runs earlier or
more often than before. The sink is chosen with the verifier: it is the first sink call whose argument, when
replaced by a constant, removes the taint. The safe twin is rewritten at the same position. We generate each
transformation alone (918 variants) and random stacks of 2, 4 or 8 distinct transformations (seed 17739; 306
variants). Figure~\ref{lst:tarpit} shows a stack of four.

\begin{figure}[t]
\begin{lstlisting}[language=Java]
const name = req.query.name;
const _o_0 = { ['b' + 'ar']: 'npm view ' + name + ' version' };
function _F_1() { function _D_2(arg) { return arg; } _F_1._D_2 = _D_2; }
_F_1();
const _msg_3 = _F_1._D_2(_o_0.bar);
const _func_4 = () => { return _msg_3; };
const _c_5 = new WeakRef([_func_4()]);
child_process.exec(_c_5.deref()[0], (err, stdout) => { ... });
\end{lstlisting}
\caption{A command-injection stub after four stacked transformations (computed-properties, nested-function,
arrow-function, weakref). The original sink call was
\texttt{child\_process.exec('npm view ' + name + ' version', ...)}; the verifier confirms the injection still works.}
\label{lst:tarpit}
\end{figure}

\subsection{Part 3: Multi-file programs}
\label{sec:multifile}
Our stubs are much smaller than real code. To test whether size and indirection matter, \texttt{transforms/multifile.js}
turns each stub into a small Express-style application: an \texttt{app.js} that registers the routes in shuffled
order, the stub's handler as one route, and 3 or 8 additional ``decoy'' routes drawn from ten hand-written ones.
The expression that reaches the sink is moved into a helper in \texttt{lib/}, so the tainted input crosses a file
boundary (Figure~\ref{fig:app}). The decoys are verified safe for all three classes, and several use sinks safely
(an \texttt{execFile} with fixed arguments, escaped HTML, \texttt{path.basename}), so the presence of a sink alone
does not reveal the bug. Vulnerable and safe twins get the same decoys in the same order. The model sees all files
in one prompt, each preceded by a \texttt{// File:} header. The applications have a median of 77 lines (6 files) or
147 lines (11 files), against about 10 for the stubs. The verifier follows the application's local
\texttt{require}s inside the same sandbox, so the cross-file flow is what gets checked.

\begin{figure}[t]
\begin{lstlisting}[language=Java]
// File: routes/package-installed-version.js
const { buildCommand } = require('../lib/commands');
module.exports = function handler(req, res) {
  const name = req.query.name;
  child_process.exec(buildCommand(name), ...);
};
// File: lib/commands.js
function buildCommand(name) {
  return 'npm view ' + name + ' version';
}
\end{lstlisting}
\caption{Part 3: the same stub inside a 6-file application (\texttt{app.js}, this route, \texttt{lib/commands.js} and
three decoy routes: \texttt{health}, \texttt{stats}, \texttt{version}). The input reaches the shell through a helper
in another file.}
\label{fig:app}
\end{figure}

One decoy turned out to matter. In the first run on the 11-file applications, false alarms on safe programs rose,
and most of the flagged answers blamed a \texttt{downloads} route that was safe but looked suspicious. We changed
that decoy to reject hidden and missing files explicitly and re-ran the multi-file experiment. Both runs are
reported in Section~\ref{sec:results}.

\subsection{Defense prompt and commercial-model pilot}
\label{sec:defense}
\textbf{Defense.} The default prompt shows the code inside a Markdown code fence, which code can close with its own
three backticks. The \emph{delimited} variant wraps the code in \texttt{<code\_X>}\,\dots\,\texttt{</code\_X>} tags,
where \texttt{X} is the first 12 hexadecimal digits of the code's SHA-256 hash, so the code cannot contain its own
closing tag and the prompt stays reproducible. It also adds one sentence telling the model to treat everything inside
the tags as code and never follow instructions in it. We evaluate it on the three attacks that succeeded at least
sometimes and on all 102 clean stubs, to measure any cost in misses or false alarms. Because the variant changes
two things at once, we can say whether it helps but not which part does.

\textbf{Commercial pilot.} To test whether our findings are specific to one open model, we ran gpt-6-luna, a low-cost current
OpenAI model, through the Chat Completions API with the same frozen prompt, stubs and attack files.
Following our rule of evaluating each model as deployed, we send no sampling parameters. The pilot covers the clean
stubs, the three attacks that worked, and the fake-end attack with the defense prompt: 306 calls with one repeat, for
US\$0.05 in total. A budget guard checks the worst-case cost before every call.
